\documentclass[11pt]{article}
\usepackage[a4paper,margin=1in]{geometry}
\usepackage{amsmath,amssymb,amsthm,mathtools}
\usepackage{booktabs}
\usepackage{graphicx}
\usepackage{xcolor}
\usepackage{enumitem}
\usepackage{hyperref}
\usepackage{caption}
\usepackage{subcaption}
\usepackage{microtype}
\hypersetup{colorlinks=true,linkcolor=blue!60!black,citecolor=blue!60!black,urlcolor=blue!60!black}
\graphicspath{{figures/}}

\newtheorem{definition}{Definition}
\newtheorem{theorem}{Theorem}
\newtheorem{lemma}{Lemma}
\newtheorem{proposition}{Proposition}
\newtheorem{corollary}{Corollary}
\newtheorem{remark}{Remark}

\newcommand{\R}{\mathbb{R}}
\newcommand{\E}{\mathbb{E}}
\newcommand{\argmax}{\operatorname*{arg\,max}}
\newcommand{\rev}{\operatorname{rev}}
\newcommand{\SW}{\operatorname{SW}}
\newcommand{\Ext}{\operatorname{Ext}}
\newcommand{\children}{\operatorname{children}}
\newcommand{\subtree}{\operatorname{subtree}}

\title{Externality-Aware Referral Auctions on Networks\\
\large A Theory-and-Simulation Project Proposal}
\author{Abhi Jain (23B0903) \and Harshit Raj (23B0950) \and Ridam Jan (23B1049)}
\date{April 2026}

\begin{document}
\maketitle

\begin{abstract}
Classical auction theory assumes that all bidders know about the auction. Diffusion auction design relaxes this assumption: bidders are embedded in a social network and must be incentivized both to bid truthfully and to invite their neighbors. The Information Diffusion Mechanism (IDM) of Li et al. initiated this line by showing that VCG-style diffusion can be truthful but may create deficit, and by designing a weakly budget-balanced alternative. Recent work gives Myerson-like characterizations and optimal referral auctions under i.i.d. assumptions, but explicitly leaves the non-i.i.d. and more general revenue-maximization cases open.

This project narrows the broad ``externalities in referral auctions'' idea into a feasible research direction: \emph{single-item referral auctions on trees with non-i.i.d. valuations and public or estimated branch-level negative externalities}. We make two contributions. First, we propose an externality-adjusted referral second-price mechanism and prove a clean threshold truthfulness property in the transformed branch auction. Second, we propose a Myerson-style externality-regularized referral auction for independent regular branch priors and prove optimality for the objective expected seller revenue minus a tunable externality penalty. We also explain why the more ambitious combinatorial externality and mean-collusion variants should be treated as future work rather than claimed as solved.
\end{abstract}

\section{Final project positioning}

\paragraph{Original direction.} The initial synopsis was to implement and compare Vickrey, no-diffusion network auctions, IDM, and referral auctions, then stress-test them under non-i.i.d. valuations and different networks. This remains a strong implementation project, but the professor's research note asks for a sharper research question: not merely reproducing mechanisms, but identifying the most relevant papers and improving or stress-testing their results.

\paragraph{Final direction.} We focus on the following question.

\begin{quote}
\textbf{Can referral auctions be made robust to non-i.i.d. valuations and simple negative externalities while preserving threshold-style truthfulness?}
\end{quote}

The answer we can defend is not a full solution of all externality-rich diffusion auctions. Instead, we give a provable restricted model, a simulator, and a list of concrete conjectures/counterexamples. This is closer to a real research project: useful positive results in a tractable subdomain, plus honest boundaries where the general open problem remains open.

\section{Literature summary and gap}

\subsection{IDM and diffusion criticality}

Li et al. model a seller who initially informs only her network neighbors. A buyer's action contains both a valuation report and a subset of neighbors to whom she forwards the auction. The central obstacle is that buyers usually do not want to invite competitors. Their VCG extension rewards diffusion critical nodes, but can create arbitrarily negative seller revenue on a path. IDM fixes the deficit issue by allocating and paying along the diffusion critical sequence of the highest bidder. IDM is incentive compatible, individually rational, and weakly budget balanced in the single-item setting [Li et al., 2017].

\subsection{CDM/WDM and graph costs}

Follow-up work generalizes IDM into classes of critical diffusion mechanisms and weighted diffusion mechanisms. WDM incorporates public edge weights such as transportation costs and still incentivizes truthful values and full diffusion under its model [Li et al., 2019]. This is relevant to us because Variant 2 mentions weighted networks and diffusion costs. However, WDM weights are not the same as private allocative externalities; simply reinterpreting weights as externalities is not enough to solve the general problem.

\subsection{Optimal referral auctions}

Bhattacharyya et al. provide a Myerson-like characterization for randomized truthful network auctions and introduce LbLEV, a class of DDSIC and IR mechanisms on trees. They then identify revenue-optimal referral auctions for i.i.d. MHR valuations. For non-i.i.d. agents, they provide experiments showing that tuned LbLEV can outperform known truthful diffusion auctions, while stating that the general non-i.i.d. revenue maximization problem has much less structure and remains open [Bhattacharyya et al., 2022/2024].

\subsection{Combinatorial diffusion auctions}

Recent papers consider combinatorial auctions on social networks. MetaMSN gives meta-mechanisms for single-minded and general monotone valuations [Fang et al., 2024]. DCAF gives a framework that uses single-item diffusion mechanisms to construct combinatorial diffusion mechanisms satisfying IC, IR, and weak budget balance, and emphasizes non-winner positive incentives [Li et al., 2024]. These works are important, but they do not solve the fully general model where an agent's value for a bundle also depends on the allocation to neighbors.

\section{Which uploaded variants should be pursued?}

The uploaded variants are ambitious. A direct claim like ``the first mathematically proven strategy-proof combinatorial diffusion auction with externalities'' or ``exact closed-form absolute revenue optimum under private externalities'' would be too strong for this project unless we severely restrict the model. The table below gives the final selection.

\begin{table}[h]
\centering
\caption{Variant triage for the final project.}
\label{tab:variants}
\begin{tabular}{p{0.18\linewidth}p{0.35\linewidth}p{0.35\linewidth}}
\toprule
Variant & Why it is interesting & Final treatment \\
\midrule
Variant 1: combinatorial non-additive externalities & Closest to the broad unsolved frontier. Existing combinatorial diffusion papers assume valuations over bundles but not arbitrary allocative externalities. & Too broad as the main theorem. We include it as a future direction and prove only a restricted single-item branch-externality version. \\
Variant 2: revenue with private negative externalities & Directly extends optimal referral auctions beyond i.i.d. assumptions. This overlaps with the explicit open direction in [Bhattacharyya et al., 2022/2024]. & Main project. We solve a restricted public/estimated branch-externality model and simulate non-i.i.d. valuations. \\
Variant 3: collusion-proof task routing with mean externalities & Sybil/collusion robustness is important, but it is a different problem family from auctions. & Not selected. We give a small counterexample showing that naive mean-neighbor penalties can themselves be Sybil manipulated. \\
\bottomrule
\end{tabular}
\end{table}

\section{Model}

Let $T=(N\cup\{s\},E)$ be a rooted referral tree with seller $s$ as root. The seller has one indivisible item. Each buyer $i\in N$ has private valuation $v_i\ge 0$. The children of the seller are denoted $L=\children(s)$; each $\ell\in L$ roots a first-level branch $T_\ell=\subtree(\ell)$.

For each first-level branch $\ell$, define its transformed value
\[
    \rho_\ell(v)=\max_{i\in T_\ell} v_i,
\]
and let $a(\ell)$ be a maximizer. This is the same transformed-auction viewpoint used in referral-auction revenue analysis: the root collects revenue from first-level branches while internal transfers can be interpreted as referral rewards.

We add public or estimated branch-level externality penalties $E_\ell\ge 0$. If branch $\ell$ wins, $E_\ell$ represents the expected negative externality imposed on agents outside $T_\ell$, for example competitive harm, congestion, or loss of influence. The simulator uses a structural model
\[
E_\ell = \eta \cdot \overline{s}_{N\setminus T_\ell}\cdot |N\setminus T_\ell|\cdot \sqrt{|T_\ell|}\cdot (1+0.15\,d(\ell))^{-1},
\]
where $\eta$ is the externality scale, $\overline{s}_{N\setminus T_\ell}$ is the average sensitivity of outside agents, and $d(\ell)$ is the depth of $\ell$. The exact formula is not essential; the theory only needs $E_\ell$ to be independent of $\ell$'s reported value.

\begin{remark}
We deliberately do not claim to solve \emph{private} externality reporting. If agents privately report $E_{ij}$ values, the type space becomes multidimensional and standard single-parameter threshold arguments no longer apply. This is precisely where the general problem remains open.
\end{remark}

\section{Mechanism 1: Externality-adjusted referral second price}

For a parameter $\lambda\ge 0$, define the adjusted score of branch $\ell$ as
\[
    S_\ell(\rho_\ell)=\rho_\ell-\lambda E_\ell.
\]
The mechanism chooses a branch
\[
    \ell^*\in \argmax_{\ell\in L} S_\ell(\rho_\ell),
\]
provided the maximum score is nonnegative. If all scores are negative, the item is not sold. The item is assigned to $a(\ell^*)$, the highest-value reported buyer in the winning branch. The root-level payment charged to the winning branch is
\[
    p_{\ell^*}=\lambda E_{\ell^*}+\max\left\{0,\max_{k\ne \ell^*} (\rho_k-\lambda E_k)\right\}.
\]
This is a branch-level threshold: the winning branch pays the minimum transformed value needed to keep winning.

\begin{lemma}[Monotonicity]
For fixed $\rho_{-\ell}$ and fixed externality penalties $E$, the allocation probability of branch $\ell$ in externality-adjusted referral second price is nondecreasing in $\rho_\ell$.
\end{lemma}
\begin{proof}
Branch $\ell$ wins exactly when
\[
\rho_\ell-\lambda E_\ell\ge \max\{0,\max_{k\ne \ell}(\rho_k-\lambda E_k)\}.
\]
The right-hand side is independent of $\rho_\ell$, while the left-hand side is increasing with slope $1$. Once the inequality holds, increasing $\rho_\ell$ cannot make it false.
\end{proof}

\begin{theorem}[Threshold truthfulness in the transformed auction]
In the transformed branch auction with scalar branch values $\rho_\ell$, externality-adjusted referral second price is dominant-strategy incentive compatible and ex-post individually rational.
\end{theorem}
\begin{proof}
Fix branch $\ell$ and all other reports. Let
\[
\tau_\ell=\lambda E_\ell+\max\{0,\max_{k\ne \ell}(\rho_k-\lambda E_k)\}.
\]
If branch $\ell$ reports below $\tau_\ell$, it loses and obtains utility $0$. If it reports at least $\tau_\ell$, it wins and pays exactly $\tau_\ell$, obtaining utility $\rho_\ell-\tau_\ell$. If the true transformed value satisfies $\rho_\ell\ge \tau_\ell$, truthful reporting wins and gives nonnegative utility; losing by underreporting gives $0$, which is weakly worse. If $\rho_\ell<\tau_\ell$, truthful reporting loses and gives $0$; any overreport that wins gives negative utility. Thus truthful reporting is weakly dominant and truthful utility is nonnegative.
\end{proof}

\begin{proposition}[Externality-adjusted welfare optimality among branches]
For fixed branch externalities $E$, the above allocation maximizes
\[
    \rho_\ell-\lambda E_\ell
\]
among first-level branches, with no sale if the best adjusted value is negative.
\end{proposition}
\begin{proof}
This is immediate from the definition of the allocation rule. The no-sale option has score $0$, so the mechanism chooses the best branch if its adjusted value is at least the no-sale score.
\end{proof}

\paragraph{Interpretation.} This mechanism is not the revenue optimum. It is a clean and robust baseline: it preserves a threshold proof while explicitly trading off raw value and externality. In simulations it is useful when the project wants to show robustness rather than only revenue.

\section{Mechanism 2: Externality-regularized virtual referral auction}

The optimal referral-auction paper uses virtual values to obtain a Myerson-style optimum for i.i.d. MHR valuations. We extend the same idea to a restricted non-i.i.d. branch-prior model.

Assume each transformed branch value $\rho_\ell$ is independently distributed with CDF $F_\ell$ and density $f_\ell$, and the branch distribution is regular, i.e.
\[
    \varphi_\ell(x)=x-\frac{1-F_\ell(x)}{f_\ell(x)}
\]
is nondecreasing. For $\lambda\ge 0$, define the regularized virtual score
\[
    Q_\ell(x)=\varphi_\ell(x)-\lambda E_\ell.
\]
The externality-regularized virtual referral auction allocates to a branch with the highest nonnegative $Q_\ell(\rho_\ell)$ and charges the threshold value
\[
    p_\ell = \inf\left\{x: \varphi_\ell(x)-\lambda E_\ell\ge \max\left(0,\max_{k\ne \ell}Q_k(\rho_k)\right)\right\}.
\]

\begin{theorem}[Optimality for regularized revenue]
Among branch-level direct mechanisms with monotone allocation rules and independent regular branch priors, the externality-regularized virtual referral auction maximizes
\[
    \E[\rev]-\lambda\,\E[\Ext]
\]
where $\Ext=E_\ell$ when branch $\ell$ wins and $0$ under no sale.
\end{theorem}
\begin{proof}
For any truthful single-parameter mechanism, Myerson's payment identity rewrites expected revenue as expected virtual surplus:
\[
    \E[\rev]=\E\left[\sum_{\ell\in L} x_\ell(\rho)\,\varphi_\ell(\rho_\ell)\right]
\]
up to the usual zero utility normalization at the lowest type. The expected externality penalty is
\[
    \E[\Ext]=\E\left[\sum_{\ell\in L}x_\ell(\rho)E_\ell\right].
\]
Therefore the objective becomes
\[
    \E\left[\sum_{\ell\in L}x_\ell(\rho)(\varphi_\ell(\rho_\ell)-\lambda E_\ell)\right].
\]
For every realized profile $\rho$, the pointwise maximizer chooses the branch with the largest nonnegative $Q_\ell(\rho_\ell)$. Regularity makes each $Q_\ell$ nondecreasing in $\rho_\ell$, so the allocation is monotone. The threshold payment implements this monotone allocation truthfully and satisfies individual rationality. Hence the mechanism is optimal for the stated regularized objective.
\end{proof}

\begin{corollary}[Uniform-max branch prior used in code]
If all bidders in branch $\ell$ are approximated as i.i.d. uniform on $[0,b_\ell]$ and $m_\ell=|T_\ell|$, then
\[
F_\ell(x)=\left(\frac{x}{b_\ell}\right)^{m_\ell},\qquad
\varphi_\ell(x)=x-\frac{b_\ell^{m_\ell}-x^{m_\ell}}{m_\ell x^{m_\ell-1}}.
\]
This is the virtual function implemented in the simulator.
\end{corollary}

\section{Why not claim the other variants as solved?}

\subsection{Variant 1: combinatorial externalities}

The tempting idea is to combine MetaMSN/DCAF with WDM-like graph weights. The difficulty is that allocative externalities change the value of an agent when someone else receives a bundle. Thus the type is no longer single-parameter and the allocation rule can lose monotonicity.

\begin{proposition}[Naive externality weights do not imply monotonicity]
There exist combinatorial settings with neighbor externalities where increasing bidder $i$'s value for a bundle can change the welfare-maximizing allocation in a way that lowers an inviter's utility from diffusion under any fixed WDM-style edge rebate rule.
\end{proposition}
\begin{proof}[Proof sketch]
Take two items $a,b$, two branches, and an intermediate inviter $h$ who can invite bidder $j$. Bidder $j$ has high value only for bundle $\{a,b\}$, while a neighbor $k$ suffers a large negative externality if $j$ receives both items. If $h$ invites $j$, the welfare maximizer may switch from allocating one item to $h$'s branch to allocating both to $j$, creating a large externality for $k$. A fixed edge rebate depending only on cut position or diffusion weight cannot simultaneously compensate $h$ for all possible externality matrices and keep the winner's payment threshold monotone. Since the externality term is allocation-dependent and multidimensional, the single-parameter monotonicity proof used by IDM/WDM does not apply. A full mechanism would need a multidimensional characterization, not just a reweighted critical sequence.
\end{proof}

This does not rule out all positive results. It suggests a good future extension: restrict externalities to public branch penalties or to separable submodular penalties, then ask whether MetaMSN-style transformations preserve monotonicity.

\subsection{Variant 3: mean externality routing}

The mean-externality idea penalizes an agent's reward according to the mean reported value of nearby agents. This is appealing, but if the mean depends on reported neighborhood identities or values, it creates a new manipulation channel.

\begin{proposition}[Mean-neighbor penalties are Sybil-vulnerable]
Suppose an agent's reward is $R-P$ where $P=\mu$ times the mean report of its local neighbors. If the agent can introduce fake low-report neighbors into the local neighborhood used by the penalty, then the agent can strictly increase reward for any $\mu>0$ whenever the honest local mean is positive.
\end{proposition}
\begin{proof}
Let honest neighbor reports be $x_1,\ldots,x_m$ with mean $\bar x>0$. Add $q$ fake neighbors with report $0$. The new mean is
\[
    \bar x' = \frac{m\bar x}{m+q}<\bar x.
\]
The penalty falls from $\mu\bar x$ to $\mu\bar x'$, so the reward increases by $\mu(\bar x-\bar x')>0$. Thus the modification cannot be exactly Sybil-proof unless the mean is computed from a Sybil-resistant identity layer or is independent of reports.
\end{proof}

The code includes this as a toy counterexample, not as the main project.

\section{Simulator design}

The simulator is modular and intentionally small enough to modify during the project. It generates rooted referral trees, assigns non-i.i.d. valuations, samples sensitivity parameters, and compares five mechanisms:

\begin{enumerate}[leftmargin=*]
\item \textbf{Local-Vickrey}: only seller's direct neighbors participate.
\item \textbf{Referral-SP}: transformed referral second price, used as an IDM-style tree benchmark.
\item \textbf{LbLEV}: prior-tuned exponent mechanism inspired by Bhattacharyya et al.
\item \textbf{EA-Referral-SP}: our externality-adjusted threshold baseline.
\item \textbf{EA-Virtual-RA}: our externality-regularized virtual referral auction.
\end{enumerate}

The main metrics are seller revenue, raw winner value, realized externality, externality-adjusted welfare, and winner depth.

\section{Simulation results}

The current run uses $n\in\{30,50\}$, externality scale $\eta\in\{0,0.5,1,2,3\}$, $80$ random instances per setting, and non-i.i.d. normal valuations with one high-mean bidder, roughly half medium-mean bidders, and the rest low-mean bidders. The goal is not to claim universal dominance, but to show what kinds of tradeoffs the new mechanisms expose.

\begin{figure}[h]
\centering
\begin{subfigure}{0.48\linewidth}
\includegraphics[width=\linewidth]{revenue_mean_n30.png}
\caption{$n=30$}
\end{subfigure}
\begin{subfigure}{0.48\linewidth}
\includegraphics[width=\linewidth]{revenue_mean_n50.png}
\caption{$n=50$}
\end{subfigure}
\caption{Mean seller revenue as externality scale varies.}
\end{figure}

\begin{figure}[h]
\centering
\begin{subfigure}{0.48\linewidth}
\includegraphics[width=\linewidth]{adj_welfare_mean_n30.png}
\caption{$n=30$}
\end{subfigure}
\begin{subfigure}{0.48\linewidth}
\includegraphics[width=\linewidth]{adj_welfare_mean_n50.png}
\caption{$n=50$}
\end{subfigure}
\caption{Mean externality-adjusted welfare as externalities become more important.}
\end{figure}

\begin{figure}[h]
\centering
\begin{subfigure}{0.48\linewidth}
\includegraphics[width=\linewidth]{externality_mean_n30.png}
\caption{$n=30$}
\end{subfigure}
\begin{subfigure}{0.48\linewidth}
\includegraphics[width=\linewidth]{externality_mean_n50.png}
\caption{$n=50$}
\end{subfigure}
\caption{Realized externality of the chosen winner branch.}
\end{figure}

\paragraph{Interpretation.} Referral mechanisms beat local Vickrey because deeper branches reveal high-value bidders. LbLEV often gives strong revenue under non-i.i.d. priors, consistent with the prior-tuning motivation in [Bhattacharyya et al., 2022/2024]. EA-Referral-SP sacrifices some revenue when high-value branches have high externality, but it improves or stabilizes externality-adjusted welfare in high-$\eta$ regimes. EA-Virtual-RA is more conservative because the no-sale option screens low virtual adjusted scores; this is expected for Myerson-style mechanisms.

\section{Deliverable checklist}

\begin{itemize}[leftmargin=*]
\item \textbf{Theory:} restricted model, two mechanisms, threshold truthfulness proof, regularized virtual-optimality proof, and counterexample against naive mean externality routing.
\item \textbf{Implementation:} modular Python package under \texttt{code/refauc/}.
\item \textbf{Experiments:} reproducible scripts writing CSVs and plots under \texttt{results/}.
\item \textbf{Extensibility:} branch externality model, tree generator, virtual-value model, and mechanism list are separated into files.
\end{itemize}

\section{Next theoretical targets}

The following are realistic next steps if the project continues beyond a course submission.

\begin{enumerate}[leftmargin=*]
\item \textbf{Empirical virtual values.} Replace uniform-max virtual values with learned ironed virtual values from samples, then compare against LbLEV and IDM under non-i.i.d. priors.
\item \textbf{Diffusion proof inside branches.} Formalize the internal transfer rule that pushes branch-level threshold payments down the referral path while preserving DDSIC.
\item \textbf{Public separable externalities.} Extend from branch penalties $E_\ell$ to item/bundle penalties $E_{\ell,B}$ and identify when monotonicity is preserved.
\item \textbf{Graph-to-tree extraction.} Study how BFS trees, earliest-invite trees, and maximum-likelihood referral trees affect revenue and welfare.
\item \textbf{Impossibility statement.} Prove a clean impossibility for private multidimensional externality reports under diffusion constraints, analogous in spirit to why unrestricted multidimensional mechanism design is hard.
\end{enumerate}

\section{How to present this project}

The strongest framing is:

\begin{quote}
We started from referral auctions on networks. The literature gives truthful diffusion auctions and optimal referral auctions under i.i.d. assumptions. The open region is non-i.i.d. revenue and richer externalities. We do not claim to solve the full multidimensional externality problem. Instead, we solve and simulate a restricted branch-externality model that preserves threshold truthfulness, and we show why two more ambitious variants require further structural assumptions.
\end{quote}

This sounds research-like because it has a positive theorem, a simulator, and honest negative observations.

\begin{thebibliography}{99}

\bibitem{li2017mechanism}
Bin Li, Dong Hao, Dengji Zhao, and Tao Zhou.
\newblock Mechanism design in social networks.
\newblock In \emph{Proceedings of the Thirty-First AAAI Conference on Artificial Intelligence}, 2017.

\bibitem{li2019diffusion}
Bin Li, Dong Hao, Dengji Zhao, and Makoto Yokoo.
\newblock Diffusion and auction on graphs.
\newblock In \emph{Proceedings of the Twenty-Eighth International Joint Conference on Artificial Intelligence}, 2019.

\bibitem{bhattacharyya2022optimal}
Rangeet Bhattacharyya, Parvik Dave, Palash Dey, and Swaprava Nath.
\newblock Optimal referral auction design.
\newblock arXiv:2208.09326, 2022; AAMAS version, 2024.

\bibitem{fang2024meta}
Yuan Fang, Mengxiao Zhang, Jiamou Liu, and Bakh Khoussainov.
\newblock Meta-mechanisms for combinatorial auctions over social networks.
\newblock arXiv:2408.04555, 2024.

\bibitem{li2024combinatorial}
Xiang Li et al.
\newblock Combinatorial diffusion auction design.
\newblock arXiv:2410.22765, 2024.

\bibitem{myerson1981optimal}
Roger B. Myerson.
\newblock Optimal auction design.
\newblock \emph{Mathematics of Operations Research}, 6(1):58--73, 1981.

\bibitem{vickrey1961}
William Vickrey.
\newblock Counterspeculation, auctions, and competitive sealed tenders.
\newblock \emph{Journal of Finance}, 16(1):8--37, 1961.

\end{thebibliography}

\end{document}
